\documentclass[final]{beamer}
\usepackage[orientation=landscape,size=custom,width=121.92,height=91.44,scale=1.15]{beamerposter}
\usetheme{I6pd2}

\usepackage[utf8]{inputenc}
\usepackage[T1]{fontenc}
\usepackage{lmodern}
\usepackage{amsmath,amssymb}
\usepackage{graphicx}
\usepackage{array}
\usepackage{booktabs}
\usepackage{tikz}
\usetikzlibrary{shapes.geometric, arrows.meta, positioning, calc}
\usepackage{xcolor}
\usepackage{enumitem}
\setlist[itemize]{leftmargin=1.2em, itemsep=0.4ex, topsep=0.3ex}

% ── Custom colors ─────────────────────────────────────────────────────
\definecolor{palmblue}{HTML}{1D4ED8}
\definecolor{palmdark}{HTML}{0F172A}
\definecolor{palmlight}{HTML}{EFF6FF}
\definecolor{palmgreen}{HTML}{16A34A}
\definecolor{palmred}{HTML}{DC2626}
\definecolor{palmorange}{HTML}{D97706}
\definecolor{palmgray}{HTML}{475569}

% Override theme colors to PALM brand
\setbeamercolor{headline}{fg=white,bg=palmdark}
\setbeamercolor{title in headline}{fg=white}
\setbeamercolor{author in headline}{fg=white}
\setbeamercolor{institute in headline}{fg=white}
\setbeamercolor{separation line}{bg=palmblue}
\setbeamercolor{block title}{fg=white,bg=palmblue}
\setbeamercolor{block body}{fg=palmdark,bg=white}

\title{PALM: Physics-Aware Leakage Minimizer}
\author{}  % Add authors here
\institute{}  % Add institution here

\begin{document}
\begin{frame}[t]

\begin{columns}[T]

% ────────────────────── COLUMN 1 ──────────────────────
\begin{column}{0.32\textwidth}

\begin{block}{Motivation}
ML models in drug discovery, materials science, and bioinformatics are routinely evaluated with \textbf{random train/test splits}. This introduces \textbf{data leakage}: structurally similar entities appear in both sets.

\vspace{0.5ex}
\begin{center}
\begin{tikzpicture}[
    box/.style={draw, rounded corners=4pt, minimum width=5cm, minimum height=1.5cm, align=center, line width=1.2pt, font=\normalsize},
]
  \node[box, fill=palmlight, draw=palmblue] (train) {\textbf{Train Set}\\[1pt]{\small Drug A, A$'$, B}};
  \node[box, fill=red!6, draw=palmred, right=2cm of train] (test) {\textbf{Test Set}\\[1pt]{\small Drug A$''$, C}};
  \draw[-{Stealth[length=4mm]}, line width=2pt, palmred, dashed]
    (train.east) -- node[above, font=\large\bfseries, palmred] {Leakage!} (test.west);
  \node[below=0.4cm of $(train)!0.5!(test)$, font=\small\color{palmgray}, align=center]
    {A, A$'$, A$''$: $>$90\% structural similarity};
\end{tikzpicture}
\end{center}

\vspace{0.5ex}
\begin{itemize}
  \item Models memorize patterns instead of learning generalizable features
  \item Reported metrics \textbf{overestimate} real-world performance
  \item Random splitting inflates test R$^2$ by \textbf{0.23--0.33} on molecular benchmarks
\end{itemize}

\vspace{0.5ex}
\textbf{PALM} provides automated, leakage-minimized splitting across molecules, materials, biomolecules, and genes --- via both CLI and web application.
\end{block}

\begin{block}{Splitting Strategies: Entity-Based \& Cluster-Based}

DataSAIL provides two complementary approaches:

\vspace{0.3ex}
{\small
\renewcommand{\arraystretch}{1.4}
\begin{tabular}{@{} >{\bfseries}l p{3.5cm} p{5.5cm} @{}}
\toprule
 & \textbf{Strategy} & \textbf{What it prevents} \\
\midrule
\color{palmred}R & Random & Nothing (baseline) \\
\color{palmorange}I1e/I1f & Entity-based (1D) & Same entity in both splits \\
\color{palmorange}I2 & Entity-based (2D) & Shared entities on either axis \\
\color{palmblue}Scaffold & Scaffold-based & Same molecular scaffold across splits \\
\color{palmgreen}C1e/C1f & Cluster-based (1D) & \textit{Similar} entities across splits \\
\color{palmgreen}C2 & Cluster-based (2D) & \textit{Similar} on both axes \\
\bottomrule
\end{tabular}
}

\vspace{0.6ex}
{\color{palmorange}\textbf{Identity (I*)}} removes \textit{exact duplicates} --- fast, no feature computation needed.

\vspace{0.3ex}
{\color{palmgreen}\textbf{Cluster (C*)}} groups \textit{similar} entities --- prevents subtle near-duplicate leakage via feature-based distance.

\vspace{0.3ex}
{\color{palmblue}\textbf{Scaffold}} groups molecules by Bemis-Murcko core ring framework --- chemistry-aware, no pairwise distance needed.

\vspace{0.3ex}
Both are essential: \textbf{I*} catches obvious leaks, \textbf{C*} catches structural analogues.
\end{block}

\begin{block}{Entity Types \& Features}
\vspace{0.2ex}
{\small
\renewcommand{\arraystretch}{1.4}
\begin{tabular}{@{} >{\bfseries\color{palmblue}}l p{9.5cm} @{}}
\toprule
Type & Feature Sets \\
\midrule
Molecule & RDKit descriptors (9), physicochemical (11), composition, Morgan FP (2048-bit ECFP4) \\
Biomolecule & Sequence properties (16), ESM2 embeddings (320--5120 dim), precomputed \\
Material & MAGPIE (120+), stoichiometry, electronic, bonding, thermodynamic, classification, mat2vec (200), CrystalNN, SOAP \\
Gene & Nucleotide composition (11), $k$-mer frequencies (80), NT embeddings (1024--2560), precomputed \\
\bottomrule
\end{tabular}
}

\vspace{0.5ex}
\textbf{Adaptive distance:} Tanimoto for fingerprints, cosine for sparse ($>$50\% zeros), PCA + Euclidean for very sparse ($>$90\%).
\end{block}

\end{column}

% ────────────────────── COLUMN 2 ──────────────────────
\begin{column}{0.32\textwidth}

\begin{block}{Results --- 1D Splitting}

\textbf{BBBP} (2,050 molecules, Morgan FP + Tanimoto):

\vspace{0.3ex}
{\small
\renewcommand{\arraystretch}{1.4}
\begin{tabular}{@{} l r r r @{}}
\toprule
\textbf{Technique} & \textbf{Mean NN} & \textbf{Zero-Dist} & \textbf{Shift} \\
\midrule
\color{palmred}Random (R) & 12.3 & 14 & 0.042 \\
\color{palmorange}Identity (I1e) & 11.6 & 8 & 0.075 \\
\color{palmgreen}\textbf{Cluster (C1e)} & \textbf{54.8} & \textbf{0} & 0.112 \\
\bottomrule
\end{tabular}
}

\vspace{0.5ex}
Identity reduces duplicate pairs (14$\to$8) but barely changes separation. Cluster: \textbf{4.5$\times$} more separation, \textbf{zero} identical pairs.

\vspace{1ex}
\textbf{Materials Project} (500 crystals, MAGPIE):

\vspace{0.3ex}
{\small
\begin{tabular}{@{} l r r r @{}}
\toprule
\textbf{Technique} & \textbf{Mean NN} & \textbf{Min NN} & \textbf{Shift} \\
\midrule
\color{palmred}Random (R) & 283.2 & 67.9 & 0.071 \\
\color{palmgreen}\textbf{Cluster (C1e)} & \textbf{494.0} & \textbf{199.5} & 0.407 \\
\bottomrule
\end{tabular}
}

\vspace{0.5ex}
Cluster: \textbf{1.7$\times$} mean, \textbf{2.9$\times$} worst-case separation.
\end{block}

\begin{block}{Results --- 2D Splitting (Drug--Target)}

\textbf{Davis DTI} (4,000 interactions, 40 drugs $\times$ 100 targets):

\vspace{0.3ex}
{\small
\renewcommand{\arraystretch}{1.4}
\begin{tabular}{@{} l r r r @{}}
\toprule
\textbf{Technique} & \textbf{Drug Ovlp} & \textbf{Target Ovlp} & \textbf{Cov.} \\
\midrule
\color{palmred}Random (R) & 40/40 (100\%) & 84/84 (100\%) & 100\% \\
\color{palmorange}Identity (I2) & 0/40 (0\%) & 0/84 (0\%) & 45\% \\
\color{palmgreen}Cluster (C1e) & 0/40 (0\%) & 84/84 (100\%) & 100\% \\
\color{palmgreen}\textbf{Cluster (C2)} & \textbf{0/40 (0\%)} & \textbf{0/84 (0\%)} & 80\% \\
\bottomrule
\end{tabular}
}

\vspace{0.5ex}
Random leaks \textbf{100\%} of entities. \textbf{I2} eliminates all overlap but loses 55\% coverage. \textbf{C1e}: zero drug overlap, full coverage. \textbf{C2}: zero overlap on both axes, 80\% coverage.
\end{block}

\begin{block}{Impact on Model Performance}

Random Forest (Morgan FP, 500 trees):

\vspace{0.3ex}
{\small
\renewcommand{\arraystretch}{1.4}
\begin{tabular}{@{} l l r r r @{}}
\toprule
\textbf{Dataset} & \textbf{Split} & \textbf{R$^2$} & \textbf{RMSE} & \textbf{MAE} \\
\midrule
ESOL & \color{palmred}Random & 0.667 & 1.226 & 0.906 \\
(1,128 mol) & \color{palmgreen}Cluster & \textbf{0.333} & 1.562 & 1.284 \\
& \color{palmgray}\textit{Inflation} & \textit{+0.334} & & \\[0.2ex]
Lipophilicity & \color{palmred}Random & 0.362 & 0.928 & 0.718 \\
(1,000 mol) & \color{palmgreen}Cluster & \textbf{0.131} & 1.115 & 0.868 \\
& \color{palmgray}\textit{Inflation} & \textit{+0.230} & & \\
\bottomrule
\end{tabular}
}

\vspace{0.5ex}
Random splitting inflates R$^2$ by \textbf{+0.23 to +0.33}. Cluster splitting with Tanimoto distance ensures test molecules are structurally novel, producing metrics that reflect genuine generalization.
\end{block}

\end{column}

% ────────────────────── COLUMN 3 ──────────────────────
\begin{column}{0.32\textwidth}

\begin{block}{System Architecture}
\begin{center}
\begin{tikzpicture}[
    every node/.style={font=\normalsize},
    stage/.style={draw=palmblue, rounded corners=4pt, minimum width=10.5cm, minimum height=1.4cm, align=center, line width=1pt, fill=palmlight, font=\normalsize\bfseries},
    detail/.style={font=\small\color{palmgray}, align=center},
    arr/.style={-{Stealth[length=3mm]}, line width=1.5pt, palmblue!60},
]
  \node[stage] (load) {1. Load Data};
  \node[detail, below=0.1cm of load] (ld) {17+ formats: CSV, SDF, FASTA, PDB, CIF, \ldots};
  \node[stage, below=0.5cm of ld] (feat) {2. Featurize Entities};
  \node[detail, below=0.1cm of feat] (fd) {4 entity types, 4--10 feature sets each (cached)};
  \node[stage, below=0.5cm of fd] (dist) {3. Compute Distances};
  \node[detail, below=0.1cm of dist] (dd) {Adaptive: Tanimoto / cosine / PCA+Euclidean};
  \node[stage, below=0.5cm of dd] (split) {4. Split (DataSAIL + Scaffold)};
  \node[detail, below=0.1cm of split] (sd) {R, I1e, I1f, I2, C1e, C1f, C2, Scaffold};
  \node[stage, below=0.5cm of sd] (eval) {5. Evaluate \& Compare};
  \node[detail, below=0.1cm of eval] (ed) {Metrics, t-SNE plots, benchmark ranking};
  \node[stage, below=0.5cm of ed, draw=palmgreen, fill=green!5] (save) {6. Save Results};
  \node[detail, below=0.1cm of save] (svd) {Format-preserving + ML exports};
  \draw[arr] (load) -- (feat); \draw[arr] (feat) -- (dist);
  \draw[arr] (dist) -- (split); \draw[arr] (split) -- (eval); \draw[arr] (eval) -- (save);
\end{tikzpicture}
\end{center}
\end{block}

\begin{block}{Web Application}
Interactive 4-step wizard (FastAPI + JavaScript):

\vspace{0.3ex}
\begin{itemize}
  \item \textbf{Upload}: drag-and-drop, ZIP, folder upload with data preview
  \item \textbf{Configure}: auto-detection of entity types, guided tooltips, technique recommendations
  \item \textbf{Run}: real-time progress via Server-Sent Events
  \item \textbf{Results}: interactive metrics dashboard, t-SNE plots, per-technique download
\end{itemize}

\vspace{0.3ex}
Built-in rate limiting, upload size limits, and automatic job cleanup for production deployment.
\end{block}

\begin{block}{Key Contributions}
\begin{itemize}
  \item \textbf{Unified framework} for leakage-aware splitting across molecules, biomolecules, materials, and genes with both entity-based (I*) and cluster-based (C*) strategies
  \item \textbf{Scaffold splitting} for chemistry-aware molecule grouping by Bemis-Murcko core
  \item \textbf{Adaptive distance selection} based on feature sparsity and type
  \item \textbf{Format-preserving output} (SDF, FASTA, CIF, \ldots) with ML-ready exports (PyTorch indices, sklearn fold column)
  \item \textbf{Benchmark mode}: automatic comparison across all techniques with recommended best split
  \item \textbf{Built-in quality metrics}: NN leakage, distribution shift, coverage, entity overlap
\end{itemize}
\end{block}

\begin{block}{Limitations \& Future Work}
\begin{itemize}
  \item C2 on very sparse features may reduce coverage; adaptive relaxation partially mitigates this
  \item Scaffold splitting is molecule-specific; fold-based (proteins) and crystal-system (materials) splitting are future work
  \item Scaling to $>$100K entities requires approximate distance methods
  \item Planned: cross-validation folds, temporal splits, batch processing
\end{itemize}
\end{block}

\begin{block}{Availability}
\begin{center}
{\Large\bfseries Open source \enspace$\bullet$\enspace Python 3.12+ \enspace$\bullet$\enspace pip installable}\\[0.5ex]
CLI: \texttt{python -m PALM config.yaml} \qquad Web: \texttt{uvicorn PALM.webapp.app:app}
\end{center}
\end{block}

\end{column}

\end{columns}

\end{frame}
\end{document}
